\begin{table}[htbp]\centering
\begin{threeparttable}
\caption{Species complexes: pairs of named species whose nearest copies lie within the within-species spread on all three markers.}
\label{tab:complexes}
\footnotesize\setlength{\tabcolsep}{4pt}
\begin{tabular}{llcrrr}
\toprule
Species A & Species B & Complex & ITS & LSU & SSU\\
\midrule
\emph{Diversispora aestuarii} & \emph{Diversispora varaderana} & C2 & 1.1 & 1.09 & 0.0\\
\emph{Diversispora arenaria} & \emph{Diversispora jakucsiae} & C3 & 3.48 & 0.83 & 0.07\\
\emph{Glomus chinense} & \emph{Glomus macrocarpum} & C1 & 5.87 & 2.18 & 0.27\\
\emph{Glomus chinense} & \emph{Glomus rugosae} & C1 & 4.95 & 1.36 & 0.2\\
\emph{Glomus macrocarpum} & \emph{Glomus rugosae} & C1 & 1.38 & 1.36 & 0.07\\
\emph{Rhizophagus irregularis} & \emph{{[}Rhizoglomus{]} vesiculiferum} & C4 & 3.25 & 2.44 & 0.34\\
\bottomrule
\end{tabular}
\begin{tablenotes}\footnotesize
\item Columns give the distance (\%) between the nearest copies of the two species. The limits are the 99th
percentile of the distance between cultures of one species: ITS 6.0\%, LSU 3.0\%, SSU 0.8\%.
\end{tablenotes}
\end{threeparttable}
\end{table}
